\begin{figure}[t!]
	\declrel{Type formation}{$\hasType T\kind$}
\begin{mathpar}
  \infrule{\rulenametypeUnit}{}{\hasType\TUnit\kindt}
  
  \infrule{{\rulenametypeArrow}}{
    \hasType T \kindt
    \\ \hasType U \kindt
  }{
    \hasType{\TFun[] T U}{\kindt}}

  % \infrule{\rulenametypeInt}{}{\hasType\TInt\kindt}
  % \infrule{\rulenametypeBool}{}{\hasType\TBool\kindt}
  \infrule{\rulenametypePair}{\hasType T \kindt \\ \hasType U \kindt}{\hasType{\TPair TU}\kindt}

  \infrule{\rulenametypeVar}{\typec{\alpha}\colon\kind \in \Delta}{\hasType\alpha \kind}
  \\
  \infrule{{\rulenametypePoly}}{
    \hasType[\Delta,\typec{\alpha}\colon\kind]{\typec{T}}{\kindt}
  }{
    \hasType{\forall (\typec{\alpha}\colon\kind) \typec{T}}{\kindt}}
  
  \infrule{\rulenametypeEndW}{}{\hasType\TEndW{\kinds}}

  \infrule{\rulenametypeEndT}{}{\hasType\TEndT{\kinds}}
 
  \infrule{\rulenametypeMsg}{\hasType T \kindp \\
    \hasType S \kinds}{\hasType
    {! T.S}{\kinds} \\ \hasType {? T.S}{\kinds}}

  \infrule{\rulenametypeDualS}{\hasType S {\kinds}}{\hasType{\TDual S} {\kinds}}

  \infrule{\rulenametypeAssume}{
    \typec{\proto : \overline{\kindp} \to \kindp} \in \Delta \\
    \hasType{\overline T}{\overline {\kindp}}
  }{\hasType{\TPApp \proto {\overline{T}}}{\kindp}}
  
  \infrule{\rulenametypeMinus}{\hasType T \kindp} {\hasType {\pneg T} \kindp}
  
  \infrule{\rulenametypeSub}{\hasType{\typec T}\kind \\ \kind \le \kind'}{\hasType{\typec T}\kind'}
  % \\
  % \infrule{\rulenametypeUp}{\hasType T \kindt}{\hasType{\TUp T} \kindp}

  % \infrule{\rulenametypePlus}{\hasType T \kindp}{\hasType {\ppos T} \kindp}

  % \textcolor{red}{\infrule{\rulenametypeSemi}{
  %   \hasType T \kindp \\ \hasType U \kindp}{
  %   \hasType {T;U} \kindp
  % }}
  % \\
  % \\
  %   % \kind \in \{\kindt, \kindp\} \\
\end{mathpar}
  \caption{Type formation}
  \label{fig:type-formation}
\end{figure}

%%% Local Variables:
%%% mode: latex
%%% TeX-master: "main"
%%% End:
